\section*{Chapitre \ref{ch:m3:automated-market-makers} --- Teneurs de marché automatisés}

\begin{solution}{exo:m3:automated-market-makers:1}
$3\,000\,000 \times 9.97/(1\,000 + 9.97) = 29\,614.74$ dollars, soit une moyenne de 2\,961.47.
\end{solution}

\begin{solution}{exo:m3:automated-market-makers:2}
$2\sqrt4/(1 + 4) - 1 = -20\%$.
\end{solution}

\begin{solution}{exo:m3:automated-market-makers:3}
Si un \omterm{def:m3:blockchains-for-traders:stable}{stablecoin} subit un \omterm{def:m3:blockchains-for-traders:stable}{décrochage} (depeg), les arbitragistes le vendent dans le pool à peu près un pour un, jusqu'à ce que le pool ne contienne plus, pour l'essentiel, que ce \omterm{def:m3:blockchains-for-traders:stable}{stablecoin}~: les
fournisseurs de liquidité finissent par détenir le \omterm{def:m3:blockchains-for-traders:stable}{stablecoin} en difficulté, acheté près du pair, comme le ferait le vendeur d'un put dont le prix d'exercice est proche du pair.
\end{solution}

\begin{solution}{exo:m3:automated-market-makers:4}
$0.8^2/8 = 8\%$ de la valeur du pool par an, soit environ 2.19 points de base par jour.
\end{solution}

\begin{solution}{exo:m3:automated-market-makers:5}
LVR~: $50\text{ millions} \times 0.25/365/8 = \$4\,281$ par jour~; frais~: $20\text{ millions} \times 0.05\% = \$10\,000$. Oui~: taux de rotation de
40~\% contre une exigence de 17.1~\%.
\end{solution}

\begin{solution}{exo:m3:automated-market-makers:6}
Sous la mesure risque-neutre, le gain espéré du portefeuille de rééquilibrage est nul, si bien que l'écart espéré entre le pool et la détention
passive égale la LVR espérée. Or la \omterm{def:m3:automated-market-makers:il}{perte impermanente} comprend aussi l'écart de risque de marché entre le pool et la détention des \omterm{def:m3:blockchains-for-traders:key}{jetons} initiaux,
que la couverture élimine~; la LVR est ce qui subsiste après couverture, et elle s'accumule trajectoire par trajectoire, quel que soit le comportement du prix.
\end{solution}

\begin{solution}{exo:m3:automated-market-makers:7}
Environ 37.0, 37.0 et 37.7 points de base sur 30 jours avec 400 trajectoires, toutes les 15 minutes, toutes les heures et tous les jours. Sans frais, la LVR espérée
ne dépend pas de la fréquence à laquelle les arbitrages ont lieu~: la valeur du pool, $2L\sqrt{P}$, perd $1 - e^{-\sigma^2t/8}$ en espérance, quel que soit
le pas de temps (36.9 points de base ici). Avec frais, un arbitrage moins fréquent laisse des écarts plus grands se creuser avant tout échange, et la fréquence
joue alors un rôle.
\end{solution}

\begin{solution}{exo:m3:automated-market-makers:8}
Les frais ne représentent que la partie recettes. Le rendement intègre aussi la perte face aux arbitragistes (LVR) et l'exposition au marché des \omterm{def:m3:blockchains-for-traders:key}{jetons} détenus~;
net de LVR et une fois couvert, le rendement peut être négatif alors que les frais ont atteint 12~\%.
\end{solution}

\begin{solution}{pb:m3:automated-market-makers:1}
\textbf{1.} $0.6^2/365/8 = 1.23$ points de base par jour~; 1\,233~dollars par jour sur 10~millions de dollars. \textbf{2.} $6\text{ millions} \times 0.30\%
\times 10/200 = \$900$. \textbf{3.} 4.11~\% de la valeur du pool par jour, soit 8.2~millions de dollars pour ce pool. \textbf{4.} $(900 - 1\,233) \times 365
\approx -\$121\,500$. \textbf{5.} La LVR quadruple~: environ $-\$1.47$ million par an. \textbf{6.} À volume égal, c'est pire~: les frais tombent à
150~\$ par jour~; un niveau de frais plus bas doit attirer plusieurs fois le volume pour être rentable. \textbf{7.} Non~: tant que le prix reste dans la plage, les frais et la LVR
croissent tous deux avec la concentration. \textbf{8.} Couvre l'exposition à l'ether (positions courtes sur \omterm{def:m3:perpetual-futures:perp}{contrats perpétuels} ou à terme, dimensionnées sur la détention d'ether du pool et
rééquilibrées), en payant le financement, les frais et le coût de rééquilibrage de la couverture elle-même. \textbf{9.} 1.23~\% par jour, si le volume ne baissait pas avec la
hausse des frais. \textbf{10.} Il prélève les frais des gros swaps sans en supporter l'arbitrage, ce qui réduit le revenu des fournisseurs passifs par unité de
LVR. \textbf{11.} Classer les swaps~: ceux qui rapprochent le pool du prix de la \omterm{def:m3:centralised-exchanges:cex}{plateforme centralisée} (CEX) dans le bloc relèvent de l'arbitrage~; les
autres, issus de routeurs et de portefeuilles connus, relèvent du bruit. \textbf{12.} Valoriser les avoirs du pool au prix de la \omterm{def:m3:centralised-exchanges:cex}{plateforme centralisée} à chaque bloc et
sommer les écarts entre le portefeuille de rééquilibrage et le pool, comme le fait l'article. \textbf{13.} Des blocs plus longs laissent se creuser de plus grands désalignements
de prix avant l'arbitrage~; en présence de frais, celui-ci attend des écarts supérieurs aux frais, et la perte par transaction dépend du temps de bloc.
\textbf{14.} Peu~: sur un mois, elle est dominée par la variation du prix~; la LVR n'apparaît qu'après couverture. \textbf{15.} P\&L couvert, frais,
LVR, composition du volume, plage de prix (si concentrée) et le \omterm{def:m3:blockchains-for-traders:gas}{gas}. \textbf{16.} Les arbitragistes et les constructeurs de blocs gagnent~; les fournisseurs passifs perdent, sauf
si le volume de bruit est élevé~; les traders de bruit paient les frais. \textbf{17.} Seulement en partie~: un fournisseur passif cote une courbe fixe, ne peut pas déplacer son prix
lorsqu'il apprend, et paie les informés. \textbf{18.} Leurs frais par unité de LVR sont élevés~: les pools de \omterm{def:m3:blockchains-for-traders:stable}{stablecoins} ont un $\sigma$ très faible et un fort
volume. \textbf{19.} \emph{Résultat nommé~:} 4.11~\% de la valeur du pool par jour pour une volatilité de 60~\% et des frais de 0.30~\%~; ce pool a un taux de rotation de 3~\% et
ne l'atteint pas. \textbf{20.} Une option que tout acteur connaissant le prix peut exploiter, payée en frais.
\end{solution}

\begin{solution}{iq:m3:automated-market-makers:1}
$(x + (1 - \gamma)\Delta x)(y - \Delta y) = xy$, donc $\Delta y = y(1 - \gamma)\Delta x/(x + (1 - \gamma)\Delta x)$.
\iqlookfor{seule l'entrée nette des frais compte dans l'invariant.}
\end{solution}

\begin{solution}{iq:m3:automated-market-makers:2}
La valeur du pool, comparée à la détention des \omterm{def:m3:blockchains-for-traders:key}{jetons} initiaux après un mouvement de prix~: $2\sqrt r/(1 + r) - 1$ pour un pool sur toute la plage de prix. Son nom laisse entendre
qu'elle disparaît~; elle ne disparaît que si le prix revient, et elle mêle une exposition au marché, couvrable, au coût réel, que la LVR isole.
\iqlookfor{la formule et la distinction avec la LVR.}
\end{solution}

\begin{solution}{iq:m3:automated-market-makers:3}
La LVR instantanée vaut $\tfrac12\sigma^2P^2|x^{*\prime}(P)|$~; avec $x^* = L/\sqrt P$, elle vaut $\tfrac14L\sigma^2\sqrt P$, et rapportée à la valeur du pool
$2L\sqrt P$, elle vaut $\sigma^2/8$.
\iqlookfor{l'argument de rééquilibrage et la détention dans un \omterm{def:m3:automated-market-makers:cp}{AMM à produit constant}.}
\end{solution}

\begin{solution}{iq:m3:automated-market-makers:4}
Les \omterm{def:m3:blockchains-for-traders:contract}{smart contracts} ne peuvent pas utiliser la virgule flottante~; les prix en racine carrée rendent linéaires les mises à jour des swaps dans une plage. Les réimplémentations hors chaîne doivent reproduire
le sens d'arrondi (vers le bas pour les sorties, vers le haut pour les entrées), les formats à virgule fixe et les bornes de tick, à l'identique, sans quoi leurs cotations et leurs exécutions simulées
divergeront de la chaîne d'une unité, ce qui suffit à faire échouer une transaction ou à mal dimensionner un arbitrage.
\iqlookfor{exactitude à la dernière unité près.}
\end{solution}

\begin{solution}{iq:m3:automated-market-makers:5}
Comparer les coûts tout compris~: impact sur le prix par pool, frais, \omterm{def:m3:blockchains-for-traders:gas}{gas} et exposition aux attaques sandwich pour un swap sur pool public~; le découpage de l'agrégateur~; pour une
intention, l'exécution attendue (dynamique des enchères, compétitivité des remplisseurs) et son risque d'échec~; et l'alternative de la \omterm{def:m3:centralised-exchanges:cex}{plateforme centralisée}. Simuler chacune
de ces options à partir de l'état courant et de l'historique d'exécutions propre à la société.
\iqlookfor{coût tout compris, y compris l'exposition à la MEV.}
\end{solution}

\begin{solution}{iq:m3:automated-market-makers:6}
Mettre à jour les prix plus vite que les arbitragistes ne peuvent négocier (courbe alimentée par un \omterm{def:m3:blockchains-for-traders:contract}{oracle}), appliquer des frais dynamiques qui croissent avec la volatilité, mettre aux enchères le droit à la
première transaction de chaque bloc et reverser le produit aux fournisseurs, ou regrouper les transactions. Chacune a un coût~: risque d'\omterm{def:m3:blockchains-for-traders:contract}{oracle}, frais plus élevés pour les
traders de bruit, complexité ou latence.
\iqlookfor{attaquer le prix obsolète dont provient la LVR, et nommer le coût.}
\end{solution}
